\subsection{相对同伦}

设 X、Y 是拓扑空间，A 是 X 的一个子空间，B 是 Y 的一个子空间。

设 \(f:\mathrm{X}\to\mathrm{Y}\) 是满足 \(f(\mathrm{A})\subset\mathrm{B}\) 的连续映射。我们称 f 是偶 (X,A) 到偶 (Y,B) 上的同伦等价，

如果存在连续映射 \(g\colon Y\to X\)（满足 \(g(B)\subset A\)）、同伦 \(\sigma\colon X\times I\to X\)（在 A 上固定，把 \(\text{Id}_X\) 与 \(g\circ f\) 相联系）以及同伦 \(\tau\colon Y\times I\to Y\)（在 B 上固定，把 \(\text{Id}_Y\) 与 \(f\circ g\) 相联系）。

假设这些条件满足。那么：

\begin{enumerate}
\def\labelenumi{\alph{enumi})}
\item
  映射 f 是从 X 到 Y 的同伦等价，映射 g 是从 Y 到 X 的同伦等价，且这两个同伦等价在同伦意义下互为逆。
\item
  此时映射 g 是偶 (Y,B) 到偶 (X,A) 上的同伦等价，称为 f 在同伦意义下的一个逆。
\item
  映射 f 与 g 通过限制到子空间，诱导出从 A 到 B 上的互逆的同胚。
\end{enumerate}

若存在偶 (X,A) 到偶 (Y,B) 上的同伦等价，则称偶 (X,A) 与 (Y,B) 同伦等价。关系“X、Y 是拓扑空间，A 是 X 的子空间，B 是 Y 的子空间，且偶 (X,A) 与 (Y,B) 同伦等价”是一个等价关系。

